\section{The ETH findings as scheme and Grothendieck-group statements}
\label{sec:eth-geometric}

\subsection{Geometric restatements of CX-E01--CX-E06}

The constructions and proofs in CX-E07--CX-E08 supply the following
dictionary. The input encoding, variable-count parameter, and
algorithmic bounds remain those in Part~I.

\begin{longtable}{@{}p{0.12\linewidth}p{0.83\linewidth}@{}}
\hline
Finding & Exact geometric restatement\\
\hline
CX-E01 &
For the succinct Boolean scheme \(X_F\), satisfiability is
nonemptiness and the satisfying-assignment count is its finite
étale degree. The theorem computes that degree in the stated
subexponential time in the number of Boolean coordinates.\\[4pt]
CX-E02 &
A decision reduction preserves emptiness; a parsimonious reduction
preserves the rank of the corresponding constant finite étale
scheme. The quantitative transfer requires the stated linear bound
on Boolean coordinates. A witness bijection also induces an
isomorphism of these constant schemes.\\[4pt]
CX-E03 &
The gate scheme projects isomorphically to the scheme of accepting
input assignments: every auxiliary coordinate is a uniquely
determined Boolean polynomial in the earlier coordinates.
Thus adding gate variables preserves both nonemptiness and degree.\\[4pt]
CX-E04 &
The degree algorithm yields the cited lower bound for Boolean
circuits of a language in \(\mathrm E^{\mathrm{NP}}\).
Its circuit size counts Boolean gates; the dictionary does not
identify these circuits with a representation of an algebraic
coordinate ring or with arithmetic circuits.\\[4pt]
CX-E05 &
The edge-selection constraints define a constant finite étale
perfect-matching scheme. Its degree is the permanent of the
specified \(0\)--\(1\) matrix; the construction has \(O(r+e)\)
Boolean coordinates and a fixed-density presentation.\\[4pt]
CX-E06 &
The computed integer is a rank of a finite free algebra.
A symbolic permanent family and a graded equivariant
orbit-closure coordinate ring require further data.
The rank statement therefore supplies neither the arithmetic
circuit construction nor the Schur-multiplicity obstruction
specified by those algebraic questions.\\
\hline
\end{longtable}

\subsection{CX-E07: counting as the degree of a finite étale scheme}
\label{finding:CX-E07}

\paragraph{Finding.}
For a Boolean \(k\)-CNF \(F\), put
\[
 X_F=\operatorname{Spec}
 \frac{\mathbb Z[x_1,\ldots,x_n]}
 {\bigl(x_i^2-x_i\ (1\leq i\leq n),\;
     \prod_{\ell\in C}(1-\ell(x))\ (C\text{ a clause of }F)\bigr)},
\]
where \(\ell(x)=x_i\) for a positive literal and
\(\ell(x)=1-x_i\) for a negative literal.
There is an isomorphism
\[
 X_F\cong\bigsqcup_{\#F}\operatorname{Spec}\mathbb Z.
\]
Consequently \(p:X_F\to\operatorname{Spec}\mathbb Z\) is finite étale,
and
\[
 \operatorname{rank}_{\mathbb Z}p_*\mathcal O_{X_F}=\#F,
 \qquad
 [p_*\mathcal O_{X_F}]=\#F[\mathcal O]\quad\hbox{in }K_0(\mathbb Z).
\]
The supplied theorem computes this degree, from the displayed
bounded-locality presentation, with its stated time and space bounds.
The presentation is succinct; an explicit list of the components or
a basis of its coordinate algebra need not be produced.

\paragraph{Proof.}
The ideals \((x_i)\) and \((x_i-1)\) are comaximal in
\(\mathbb Z[x_i]\).
The Chinese remainder theorem, applied successively in all variables,
identifies
\[
 \mathbb Z[x_1,\ldots,x_n]/(x_i^2-x_i)_i
 \cong\prod_{a\in\{0,1\}^n}\mathbb Z
\]
by evaluation.
The polynomial attached to a clause evaluates to one at an assignment
violating that clause, and to zero at one satisfying it.
Quotienting by all clause polynomials therefore removes precisely
the factors corresponding to nonsatisfying assignments.
The surviving factors give the asserted disjoint-union isomorphism.
Their coordinate module is free of rank \(\#F\), which proves the
rank and \(K_0\) identities.
The input formula and the factored clause presentation use the same
literal data, so the counter computes that rank without changing the
parameters beyond fixed encoding constants.
\hfill\(\square\)

\paragraph{The exact dictionary for CX-E01--CX-E05.}
Satisfiability is nonemptiness of \(X_F\).
A witness chooses one of its components.
Counting computes the degree of \(p\).
A parsimonious reduction in CX-E02 preserves the number of components
and hence the isomorphism class of the constant finite étale scheme.
A specified witness bijection identifies the component sets.
In the gate constructions of CX-E03 and CX-E05 there is a more
explicit isomorphism: projection forgets gate variables, and its
inverse assigns each gate the multilinear Boolean polynomial of
its predecessors.
For each Boolean input, the gate clauses permit exactly that value;
the product decomposition proved above therefore shows that the
projection and its inverse compose to the identity as morphisms
over \(\mathbb Z\).
The gate variables introduce no extra components.
All quantitative bounds remain bounds in the size of the specified
presentation; they are not complexity assertions about a scheme
given only up to an unencoded isomorphism.

\paragraph{Exterior powers, weights, and the exact limits.}
In the \(\lambda\)-ring of vector bundles,
\[
 \lambda_t([p_*\mathcal O_{X_F}])=(1+t)^{\#F},
\]
because the \(j\)-th exterior power of a free module of rank \(\#F\)
has rank \(\binom{\#F}{j}\).
Its reduced rank class is zero, so this construction produces no
nonzero class in the positive \(\gamma\)-filtration.
For a fiber over \(\mathbb F_q\) and \(\ell\ne\operatorname{char}\mathbb F_q\),
the geometric \(\ell\)-adic cohomology is concentrated in degree zero,
has dimension \(\#F\), and has identity Frobenius.
The fiber's zeta function is consequently
\((1-t)^{-\#F}\).
These facts follow component by component from the proved
decomposition. The object has weight zero and its Artin motive is
a sum of unit objects; no new purity theorem is asserted.

\subsection{CX-E08: the additional data a GCT obstruction would require}
\label{finding:CX-E08}

\paragraph{Finding.}
The rank computation in CX-E07 does not, through this dictionary,
construct a separating representation in a determinant-orbit-closure
coordinate ring.

\paragraph{Exact comparison and proof.}
Let \(E=\mathbb Z^{m^2}\), let \(G=\mathrm{GL}(E)\), and let
\(\mathbb V\) be the affine coefficient space associated to
\(\operatorname{Sym}^m(E^*)\).
Take the schematic closure in \(\mathbb V\) of the determinant orbit.
After base change to the complex numbers, a containment of a padded
permanent orbit closure \(Y\) in the determinant orbit closure \(X\)
would give a graded \(G\)-equivariant quotient
\[
 \mathbb C[X]\longrightarrow\mathbb C[Y].
\]
In every finite homogeneous degree, rational representations of this
general linear group are semisimple.
Write \(E_{\mathbb C}=E\otimes_{\mathbb Z}\mathbb C\).
For each irreducible polynomial representation
\(S_\lambda E_{\mathbb C}\), equivariance
therefore sends its isotypical source summand onto the target
summand. Writing the two summands as
\(S_\lambda E_{\mathbb C}\otimes M_\lambda\) and
\(S_\lambda E_{\mathbb C}\otimes N_\lambda\), surjectivity implies
\(\dim N_\lambda\leq\dim M_\lambda\).
A reversed inequality would obstruct the containment
\cite[\S7.4, (7.6), p.~38, and (7.9), p.~39]{Burgisser2024}.
This is a statement in the Grothendieck group of rational
representations, with Schur and exterior-power operations, rather
than a statement about total rank alone.

CX-E07 provides a bounded-degree Boolean presentation and the integer
\(\#F\).
Its decomposition into unit components supplies no such graded
\(G\)-representation, map of coordinate rings, or inequality of Schur
multiplicities.
Consequently this proved dictionary does not produce that obstruction.
A new construction furnishing these additional data, with the
appropriate complexity guarantees, would be needed.
This verifies the missing data in the attempted implication; it does
not rule out a future construction.
\hfill\(\square\)

The exact EGA, SGA, FGA and archival locators for the constructions in
this section are given in the paper's common geometric source
dictionary. The mathematical identifications above use only the
displayed decomposition, finite-module operations, and the
characteristic-zero representation statement explicitly indicated.
